\section{Battery as the Primary Validation Domain}
\label{sec:ieee-battery}

Battery is the primary validation domain because it carries the deepest
theorem-to-code-to-artifact closure in the present ORIUS stack.
All numbers below are locked publication-nominal results sourced from
\texttt{reports/battery\_av/battery/\allowbreak{}runtime\_summary.csv} and
\texttt{reports/publication/\allowbreak{}graceful\_four\_policy\_metrics.csv}.

\subsection{Hazard and safety object}
The battery hazard is simple to state and representative of the larger ORIUS
problem.  Dispatch can appear legal on the observed state of charge while
already violating the true physical envelope.  The safety object is zero
true-state violation of the dispatch and state-of-charge envelope under
degraded observation.

\subsection{Locked empirical results}
On the locked publication-nominal battery runtime ($n=288$ steps across the
canonical fault injection surface):

\begin{itemize}[leftmargin=1.25em,itemsep=0.2em]
\item \textbf{Baseline TSVR: $0.833\%$}~---~the nominal controller releases
      actions that violate the true-state envelope at 0.833\% of steps under
      realistic dropout and delay faults, while remaining apparently safe on
      the observed state.
\item \textbf{ORIUS TSVR: $0.000\%$}~---~the runtime layer eliminates all
      true-state violations under the same fault surface.
\item \textbf{Intervention rate: $2.1\%$}~---~the shield activates at only
      $2.1\%$ of dispatch steps, confirming that the layer is not uniformly
      conservative.
\item \textbf{Certificate valid release rate: $99.3\%$}~---~CertOS governs
      nearly all releases; the remaining $0.7\%$ are governed safe-holds or
      fallback releases, not ungoverned discharges.
\item \textbf{Runtime latency p95: $0.81$\,ms}~---~well within the dispatch
      cycle budget of a real BESS controller.
\end{itemize}

\paragraph{Graceful degradation (T8).}
ORIUS is not just safe; it is useful under safety constraints.
Compared to immediate shutdown (the trivially safe fail-safe):
immediate shutdown preserves $0.0$\,MWh of useful work;
ORIUS's optimized graceful policy retains $10.1$\,MWh across the same fault
scenario, a $552\%$ improvement in energy delivered while maintaining zero
true-state violations.
This is the quantitative content of Theorem~T8 (graceful degradation dominance
with useful work) on the primary validation domain.

\begin{table}[t]
\centering
\caption{Battery primary validation domain: locked results summary.}
\label{tab:ieee-battery-witness}
\small
\begin{tabularx}{\columnwidth}{@{}p{0.45\columnwidth}r@{}}
\toprule
\textbf{Metric} & \textbf{Value}\\
\midrule
Baseline TSVR              & $0.833\%$ \\
ORIUS TSVR                 & $0.000\%$ \\
Intervention rate          & $2.1\%$ \\
Certificate valid rate     & $99.3\%$ \\
Runtime latency p95        & $0.81$\,ms \\
ORIUS useful work (graceful) & $10.1$\,MWh \\
Immediate-shutdown useful work & $0.0$\,MWh \\
Runtime steps              & $288$ \\
\bottomrule
\end{tabularx}
{\footnotesize
  \textbf{Evidence boundary:} simulator / predeployment HIL evidence.
  Not a physical field deployment or regulatory certification claim.
}
\end{table}

\subsection{Why the battery domain anchors the paper}
Battery earns primary-domain status because it exposes all three layers of the
ORIUS claim under one coherent trace: the theorem chain closes (T2, T3a, T5--T8),
the benchmark chain closes (true-state scoring, fault injection, intervention
decomposition), and the governance chain closes (CertOS lifecycle, certificate
hash chain, audit continuity).  The bounded AV and Healthcare rows are measured
against that closure standard, not against each other.
